% GENERATED by build_tex.py (preprint: real commit hashes)
\section{Methods and definitions}\label{app:methods}
% DRAFT-TEXT: review (one framing sentence; the table quotes the committed definitions)
Each term is given in the wording of the file that fixed it, with that file and the commit of the quoted text; code lines are quoted verbatim. \emph{Certified} means that all four criteria hold at the point: a numerically certified local minimum under these tolerances, not a global or unique one. Stop reasons, with solver errors and hangs separated from caps, are tabulated in \texttt{paper2/revision/uncertainty\_output.txt}.
{\footnotesize\setlength{\tabcolsep}{3pt}
\begin{longtable}{R{0.16\linewidth} R{0.50\linewidth} R{0.16\linewidth} R{0.08\linewidth}}
\toprule
Term & Committed definition & File & Commit \\
\midrule
\endhead
Edge function (Gaussian RBF + SiLU) & Same edge type: SiLU base plus 8 Gaussian RBFs on fixed centres linspace($-$1, 1, 8), width 2/7. & \texttt{paper2/\allowbreak{}synth\_\allowbreak{}prereg.md} & dfd281d \\
P0 (second differences) & \textbf{P0, "D2 only":} $\lambda$ $\cdot$ $\Sigma$\_e $\|$D2 c\_e$\|^{2}$. This is the final.py smoothness penalty. & \texttt{paper2/\allowbreak{}synth\_\allowbreak{}prereg.md} & dfd281d \\
P1 (adds the SiLU-weight term) & \textbf{P1, "curvature penalty, Bagrow-style":} $\lambda$ $\cdot$ $\Sigma$\_e ($\|$D2 c\_e$\|^{2}$ + K\_silu $\cdot$ $\alpha$\_e$^{2}$). & \texttt{paper2/\allowbreak{}synth\_\allowbreak{}prereg.md} & dfd281d \\
P2 (null-space penalty) & \textbf{P2, "null space closed":} P1 + $\mu$ $\cdot$ Q, with Q = $\Sigma$\_e $\|$P\_null c\_e$\|^{2}$. & \texttt{paper2/\allowbreak{}synth\_\allowbreak{}prereg.md} & dfd281d \\
Null-space penalty on real data & Loss: the final.py loss + $10^{-4}$ $\times$ Q, with Q from \texttt{close.py} (the RBF null-space part of every edge + the squared SiLU base weights). & \texttt{paper2/\allowbreak{}revision/\allowbreak{}infimum\_\allowbreak{}prereg.md} & 349022b \\
Polishing & \textbf{"polished":} trust-exact from the early state, in float64, cap 500 iterations, on the cell's loss. The callback is named \texttt{intermediate\_\allowbreak{}result}, with a 2 h safety cap per run. & \texttt{paper2/\allowbreak{}synth\_\allowbreak{}prereg.md} & dfd281d \\
Certified (definition) & \textbf{CERTIFIED} (the certify.py definition, on the cell's loss): max|grad| < $10^{-7}$; smallest eigenvalue $\geq$ $-10^{-8}$ $\times$ largest; Newton decrement on the non-null eigenspace < $10^{-16}$; $\|$g\_null$\|$ < $10^{-9}$. & \texttt{paper2/\allowbreak{}synth\_\allowbreak{}prereg.md} & dfd281d \\
Certification tolerances (code) & \texttt{GTOL, EIGREL, NEWTON, GNULL = 1e-7, 1e-8, 1e-16, 1e-9} & \texttt{certify.py}:33 & 36f295b \\
Certification criteria (code) & \texttt{crit = dict(a=gmax < GTOL, b=bool(ev[0] >= -EIGREL * lmax), c=newton < NEWTON, d=gnull < GNULL)} & \texttt{certify.py}:50 & 36f295b \\
Near-null directions (code) & \texttt{null = np.abs(ev) < EIGREL * lmax} & \texttt{certify.py}:45 & 36f295b \\
Exact symmetries (B-spline null set) & The null set is the 15-dimensional constant-shift space, constructed analytically from the parameterisation (never fitted from observed flat directions), per the shared-curve Step 0 precedent. & \texttt{paper2/\allowbreak{}revision/\allowbreak{}bspline\_\allowbreak{}prereg.md} & 03638f8 \\
Null-set-aware certification & \textbf{Certification of the $\gamma$ = 0 training objectives} (data loss + P0 or P1) is null-set-aware: $\lambda$\_min, the Newton decrement (a complement-restricted solve) and the null-space-gradient criterion are evaluated in the complement. & \texttt{paper2/\allowbreak{}revision/\allowbreak{}bspline\_\allowbreak{}prereg.md} & 03638f8 \\
Degenerate certified point & \textbf{DEGENERATE.} A certified point is DEGENERATE if it has more than 3 near-flat directions OR saturation above 10\%. (D2: committed synthetic P2 polished models have 0.) & \texttt{paper2/\allowbreak{}revision/\allowbreak{}infimum\_\allowbreak{}prereg.md} & 179814f \\
Probes at degenerate points & \textbf{Three probes at every degenerate point.} All are evaluated on L + $\gamma\|\theta\|^{2}$ at that $\gamma$. 1. \textbf{Restart:} trust-exact with the gradient tolerance tightened $\times$100 ($10^{-7}$ $\to$\allowbreak{} $10^{-9}$), with the same caps. 2. \textbf{Scaling:} multiply (i) the layer-1 affine RBF component, and (ii) all layer-1 weights and biases, by c $\in$ \{1.01, 1.1, 1.5, 2, 5, 10\}. 3. \textbf{Flat-direction search:} the up to 10 flattest Hessian eigenvectors (smallest |$\lambda$|), both signs, loss at 13 log-spaced steps s $\in$ \{10\^{}-3, 10\^{}-2.75, \ldots{}, 10\^{}0\} $\times$ $\|\theta\|$ ($\theta$ = all parameters). & \texttt{paper2/\allowbreak{}revision/\allowbreak{}infimum\_\allowbreak{}prereg.md} & 179814f \\
PLATEAU / FLAT MINIMUM & \textbf{PLATEAU} if either of these holds: & \texttt{paper2/\allowbreak{}revision/\allowbreak{}infimum\_\allowbreak{}prereg.md} & 179814f \\
Stalled restart & \textbf{Stalled restart.} & \texttt{paper2/\allowbreak{}revision/\allowbreak{}infimum\_\allowbreak{}prereg.md} & 179814f \\
Path exponent (I2) & \textbf{I2.} Runaway models: fit log $\|\theta^{*}\|^{2}$ = a $-$ $\alpha$ log $\gamma$ over the smallest four certified (incl. certified-degenerate) $\gamma$ values. $\alpha$ must exceed 0.2, with a bootstrap 95\% interval that excludes 0, in $\geq$ 80\% of models. & \texttt{paper2/\allowbreak{}revision/\allowbreak{}infimum\_\allowbreak{}prereg.md} & 179814f \\
Plateau ratio (I3) & \textbf{I3.} Control (P2) models plateau: the $\|\theta^{*}\|^{2}$ ratio between the two smallest certified $\gamma$ values is $\leq$ 1.5, in $\geq$ 80\% of models. & \texttt{paper2/\allowbreak{}revision/\allowbreak{}infimum\_\allowbreak{}prereg.md} & 179814f \\
Short-path statistic D & D = $\|\theta^{*}$($10^{-7}$)$\|^{2}$ / $\|\theta^{*}$($10^{-5}$)$\|^{2}$, computed from the points as reached (primary; this is the D that enters R2 and R4). Both points are flagged if either is not certified. A certified-only variant (models with both points certified, incl. certified-degenerate) is reported beside it. & \texttt{paper2/\allowbreak{}revision/\allowbreak{}dose\_\allowbreak{}prereg.md} & 1a18569 \\
Seed disagreement d\_AB & d\_AB between the 28 seed pairs within a cell: RMS($\Delta$pred) / std(teacher output, training rows). & \texttt{paper2/\allowbreak{}synth\_\allowbreak{}prereg.md} & dfd281d \\
Unit match & \textbf{unit match:} Hungarian on |corr| of unit contributions over the training rows, median over 3 units $\times$ 8 seeds; & \texttt{paper2/\allowbreak{}synth\_\allowbreak{}prereg.md} & dfd281d \\
Cross-edge cancellation ratio & \textbf{Logged at both states:} data term, penalty term(s), $\mu Q$, $\|\theta\|^{2}$ per block (layer-1 RBF, layer-2 RBF, SiLU base, biases), the median per-unit cross-edge cancellation ratio ($\Sigma$\_i RMS(edge\_ij) / RMS($\Sigma$\_i edge\_ij) over the 4 incoming edges, on the training rows), the saturated fraction (|s\_j| > 3), max|grad|, and an eigenvalue summary (min, max, number negative < $-10^{-8}$ $\times$ max, null-set size, deciles). & \texttt{paper2/\allowbreak{}synth\_\allowbreak{}prereg.md} & dfd281d \\
\bottomrule
\end{longtable}}
\noindent Architecture per experiment, from the code (normalisation layers searched in all 65 tracked Python files: none found):
{\footnotesize\setlength{\tabcolsep}{3pt}
\begin{longtable}{R{0.13\linewidth} R{0.10\linewidth} R{0.25\linewidth} R{0.42\linewidth}}
\toprule
Experiment & LayerNorm & Biases and interlayer nonlinearity & Code (file:line, committed) \\
\midrule
\endhead
Synthetic study (Gaussian RBF), [4, 3, 1] & none & one bias per unit in each layer (the base Linear); tanh between the layers; no salt embedding; inputs clipped to [-1, 1] & \texttt{paper2/\allowbreak{}synth.py}:55 \texttt{class SynthKAN(\allowbreak{}nn.Module):}; \texttt{kan.py}:42 \texttt{self.base = nn.Linear(\allowbreak{}d\_\allowbreak{}in, \allowbreak{}d\_\allowbreak{}out)}; \texttt{paper2/\allowbreak{}synth.py}:63 \texttt{return self.layers[1](\allowbreak{}torch.tanh(\allowbreak{}self.layers[0](\allowbreak{}x)))} \\
B-spline replication, [4, 3, 1] & none & one bias per unit in each layer; tanh between the layers & \texttt{paper2/\allowbreak{}revision/\allowbreak{}bspline\_\allowbreak{}basis.py}:76 \texttt{class BSynthKAN(\allowbreak{}nn.Module):}; \texttt{paper2/\allowbreak{}revision/\allowbreak{}bspline\_\allowbreak{}basis.py}:59 \texttt{self.base = nn.Linear(\allowbreak{}d\_\allowbreak{}in, \allowbreak{}d\_\allowbreak{}out)}; \texttt{paper2/\allowbreak{}revision/\allowbreak{}bspline\_\allowbreak{}basis.py}:84 \texttt{return self.layers[1](\allowbreak{}torch.tanh(\allowbreak{}self.layers[0](\allowbreak{}x)))} \\
Real data, KAN-$\delta$ [10, 6, 1] & none & one bias per unit in each layer; 2-dimensional salt embedding passed through tanh; tanh between the layers; output multiplied by x\_co (delta\_head) & \texttt{kan.py}:60 \texttt{class KAN(\allowbreak{}nn.Module):}; \texttt{kan.py}:42 \texttt{self.base = nn.Linear(\allowbreak{}d\_\allowbreak{}in, \allowbreak{}d\_\allowbreak{}out)}; \texttt{kan.py}:68 \texttt{z = torch.cat(\allowbreak{}[x, \allowbreak{}torch.tanh(\allowbreak{}self.emb(\allowbreak{}salt))], \allowbreak{}1)  \# tanh keeps the embedding on the RBF grid}; \texttt{kan.py}:70 \texttt{z = layer(\allowbreak{}torch.tanh(\allowbreak{}z) if i else z)  \# hidden activations squashed onto [-1, \allowbreak{}1] before the next grid} \\
Shared-curve KAN (SC-KAN-$\delta$) & none & biases in the mixing layer and the head, none in the salt map; tanh after the mixing layer; the shared curves are mean-centred (a fixed centre buffer), which is not a normalisation layer & \texttt{arch/\allowbreak{}sc\_\allowbreak{}model.py}:30 \texttt{class SCKAN(\allowbreak{}nn.Module):}; \texttt{arch/\allowbreak{}sc\_\allowbreak{}model.py}:37 \texttt{self.mix = nn.Linear(\allowbreak{}d, \allowbreak{}H)                                       \# weight = A\^{}T (\allowbreak{}H x d), \allowbreak{}bias = b}; \texttt{arch/\allowbreak{}sc\_\allowbreak{}model.py}:39 \texttt{self.salt = nn.Linear(\allowbreak{}2, \allowbreak{}H, \allowbreak{}bias=False)                          \# W (\allowbreak{}H x 2)}; \texttt{arch/\allowbreak{}sc\_\allowbreak{}model.py}:40 \texttt{self.head = nn.Linear(\allowbreak{}H, \allowbreak{}1)                                      \# w, \allowbreak{}b0}; \texttt{arch/\allowbreak{}sc\_\allowbreak{}model.py}:53 \texttt{return self.head(\allowbreak{}torch.tanh(\allowbreak{}self.hidden(\allowbreak{}x, \allowbreak{}s, \allowbreak{}fit)))[:, \allowbreak{}0]}; \texttt{arch/\allowbreak{}sc\_\allowbreak{}model.py}:49 \texttt{z = g - (\allowbreak{}g.mean(\allowbreak{}0) if fit else self.center)} \\
Two-anchor excess model (EX-KAN) & none & two KAN subnetworks (s: [4, 3, 1]; g: kan.KAN [7, 6, 1] with salt embedding), each with one bias per unit per layer and tanh between layers & \texttt{arch/\allowbreak{}ex\_\allowbreak{}model.py}:33 \texttt{class EXKAN(\allowbreak{}nn.Module):}; \texttt{arch/\allowbreak{}ex\_\allowbreak{}model.py}:24 \texttt{self.layers = nn.ModuleList(\allowbreak{}[KANLayer(\allowbreak{}4, \allowbreak{}3), \allowbreak{}KANLayer(\allowbreak{}3, \allowbreak{}1)])}; \texttt{arch/\allowbreak{}ex\_\allowbreak{}model.py}:27 \texttt{return self.layers[1](\allowbreak{}torch.tanh(\allowbreak{}self.layers[0](\allowbreak{}x)))[:, \allowbreak{}0]}; \texttt{arch/\allowbreak{}ex\_\allowbreak{}model.py}:37 \texttt{self.g = KAN(\allowbreak{}len(\allowbreak{}G\_\allowbreak{}COLS), \allowbreak{}[6, \allowbreak{}1], \allowbreak{}n\_\allowbreak{}salt)} \\
\bottomrule
\end{longtable}}